\begin{lemma}{Eigenschaften von Ordnungsstatistiken}
             {EigenschaftenVonOrdnungsstatistiken}
Seien $n \in \N$, $k \in \haken n$ und $\gamma \in \mcm_n(\Omega,\mfa,\R^n)$.\\
Dann gilt:
\begin{enumerate}[(1)]
	\item\label{EigenschaftenVonOrdnungsstatistiken1}
	$\gamma^\sigma_k, \gamma^\ast_k \in \mcm(\Omega, \mfa, \R)$.
	\item\label{EigenschaftenVonOrdnungsstatistiken2}
	Ist $\Bild \gamma \subseteq [0, \infty[^n$, so ist
	$\gamma_k^\ast = \gamma^\sigma_k$. Ist $\gamma$ unabh�ngig, so ist auch
	\[\gamma_{\abs{\cdot}}: \Omega \rightarrow \R^n, \omega \mapsto
	(\abs{\gamma_i(\omega)})_{i=1}^n\]
	unabh�ngig.
	\item\label{EigenschaftenVonOrdnungsstatistiken3}
	Ist $\gamma$ unabh�ngig und identisch verteilt,
	so gilt:
	\[\forall\, t \in \R: \Prim(\gamma^\sigma_k > t) 
	=\sum_{l=0}^{n-k}\binom n l
	\Prim(\gamma_1\le t)^l(1-\Prim(\gamma_1\le t))^{n-l}\text.\]
	\item\label{EigenschaftenVonOrdnungsstatistiken4}
	Ist $\gamma$ unabh�ngig und identisch verteilt, sodass $\gamma_1$ eine
	stetige Dichte $f$ besitzt, so ist 
	$\Prim(\gamma^\sigma_k \le \cdot)$ stetig differenzierbar und es gilt
	f�r alle $t_0 \in \R$:
	\[\left(\frac d {dt} \Prim(\gamma^\sigma_k\le \cdot) \right)(t_0) 
	= n\binom{n-1}{n-k}
	\Prim(\gamma_1 \le t_0)^{n-k}(1-\Prim(\gamma_1 \le t_0))^{k-1}f(t_0)\text.\]
\end{enumerate}
\end{lemma}
\begin{proof}
(1) Nach
\refclaim{EigenschaftenDesKtenMaximums}{EigenschaftenDesKtenMaximums2c}
ist $\mu_{k,n}$ stetig, also $\mfb_n-\mfb$--messbar und daher ist
$\gamma^\sigma_k=\mu_{k,n} \circ\gamma$ selbst $\mfa-\mfb$--messbar als
Komposition messbarer Funktionen.\\
Die Abbildung $\pfeil {\abs\cdot}n$ ist offenbar stetig und daher ist
$\gamma_k^\ast = \mu_{k,n} \circ\pfeil {\abs\cdot}n\circ \gamma$ als
Komposition stetiger und messbarer Funktionen $\mfa-\mfb$--messbar.\\
(2) Die erste Aussage ist direkt aus der Definition ersichtlich.\\
F�r die zweite seien $(A_i)_{i \in \haken n} \in \mfb^n$. Da die Abbildung
$f:\R \rightarrow \R, x \mapsto -x$ stetig und damit messbar ist, ist
$-A_i = f\inv(A_i) \in \mfb$ f�r alle $i \in \haken n$. Also gilt:
\begin{align*}
&\,\Prim\left(\bigcap_{i \in \haken n} \meng{\abs{\gamma_i} \in A_i} \right)
= \Prim\left(\bigcap_{i \in \haken n} 
\meng{\gamma_i \in A_i \vee \gamma_i \in -A_i} \right)\\
=&\,\Prim\left(\bigcap_{i\in\haken n}\meng{\gamma_i\in(A_i \cup -A_i)} \right)
\overset{\gamma \text{ unabh.}}=
\prod_{i \in \haken n} \Prim\left(\meng{\gamma_i \in (A_i \cup -A_i )}\right)\\
=&\prod_{i \in \haken n} \Prim\left(\meng{\abs{\gamma_i} \in A_i)}\right)\text.
\end{align*}
(3)/(4) Siehe \cite{AB}, dort 2.2.13 f�r die Verteilung und 2.2.2 f�r die
Dichte. Dabei ist zu beachten, dass in dem genannten Buch Ordnungsstatistiken
andersherum definiert sind, d. h. aufsteigend.
\end{proof}